% ------------------------------------------------------------------------------
% 4 Results I -- Treatment effects
% ------------------------------------------------------------------------------
\section{Results I: Treatment Effects}\label{sec:res-tx}

\subsection{Kaplan--Meier analysis}

Figure~\ref{fig:km} displays Kaplan--Meier estimates for the composite
DFS endpoint and for overall survival, each with the numbers at risk.
The three curves separate early and the separation widens over the first
three years, driven almost entirely by the levamisole-plus-5-FU arm.
The overall log-rank test is highly significant for DFS
($\chi^2 = 21.0$ on 2 degrees of freedom, $p < 0.001$); pairwise
comparisons show that levamisole alone is indistinguishable from
observation (HR $0.968$, $p = 0.76$), whereas levamisole plus 5-FU is
superior to both observation (HR $0.623$, $p < 0.001$) and to
levamisole alone ($p < 0.001$). The pattern for overall survival is
similar but attenuated (adjusted HR $0.688$, 95\% CI $0.544$--$0.871$),
consistent with post-recognition salvage therapies diluting survival
differences.

\begin{figure}[htbp]
  \centering
  \includegraphics[width=0.86\textwidth, max height=0.40\textheight]{fig_km_dfs.png}\\[4pt]
  \includegraphics[width=0.86\textwidth, max height=0.40\textheight]{fig_km_os.png}
  \caption{Kaplan--Meier estimates with pointwise 95\% confidence bands
  and numbers at risk. Top: composite disease-free survival (primary
  endpoint). Bottom: overall survival (secondary endpoint).}
  \label{fig:km}
\end{figure}

\subsection{Cox models}

Table~\ref{tab:cox} assembles the treatment-effect estimates across
specifications. Adjustment for the ten baseline covariates, pooled over
20 multiple imputations by Rubin's rules, moves the primary contrast
only marginally (unadjusted HR $0.623$ vs.\ adjusted $0.621$), as
expected under randomization with balanced arms; the complete-case
analysis (not shown in full) agrees closely. The adjusted covariate
effects themselves (Figure~\ref{fig:forest}) are clinically coherent:
each additional positive lymph node multiplies the hazard by
$1.08$ (95\% CI $1.06$--$1.10$), spread to contiguous structures
triples the hazard relative to submucosal involvement
(HR $3.15$, 95\% CI $1.30$--$7.65$), obstruction raises it by 25\%, and
a long surgery-to-registration interval is protective in this cohort
(HR $0.76$, 95\% CI $0.63$--$0.92$)---the latter reflecting the
dataset's coding of a longer recovery interval before registration
among patients who recovered well enough to be registered.

\begin{table}[htbp]
\centering\small
\caption{Treatment-effect estimates for the composite DFS endpoint
across analysis specifications. The Fine--Gray model targets the
recurrence endpoint with death without recurrence as a competing event.}
\label{tab:cox}
\begin{tabular*}{\textwidth}{@{\extracolsep{\fill}} l l c c c}
\toprule
Analysis & Contrast & HR & 95\% CI & $p$ \\
\midrule
Unadjusted Cox & Lev vs.\ Obs & 0.968 & 0.790--1.187 & 0.757 \\
Unadjusted Cox & Lev+5FU vs.\ Obs & 0.623 & 0.499--0.777 & $<0.001$ \\
Adjusted Cox, MI ($m=20$) & Lev vs.\ Obs & 0.947 & 0.770--1.165 & 0.607 \\
Adjusted Cox, MI ($m=20$) & Lev+5FU vs.\ Obs & \textbf{0.621} & \textbf{0.496--0.776} & $<0.001$ \\
Adjusted Cox, MI (OS endpoint) & Lev+5FU vs.\ Obs & 0.688 & 0.544--0.871 & 0.002 \\
Fine--Gray subdistribution & Lev+5FU vs.\ Obs & 0.596 & 0.469--0.756 & $<0.001$ \\
\bottomrule
\end{tabular*}
\end{table}

\begin{figure}[htbp]
  \centering
  \includegraphics[width=0.82\textwidth, max height=0.42\textheight]{fig_forest_adjusted.png}
  \caption{Adjusted hazard ratios for disease-free survival from the
  multiply imputed Cox model (log scale). Red indicates estimates whose
  95\% confidence interval excludes 1.}
  \label{fig:forest}
\end{figure}

\subsection{Absolute benefit}

Five-year Kaplan--Meier DFS probabilities were 42.4\% (95\% CI
36.9--47.8) under observation, 44.2\% (38.6--49.6) under levamisole,
and 59.2\% (53.4--64.5) under levamisole plus 5-FU; the corresponding
overall-survival probabilities were 52.6\%, 53.5\%, and 63.4\%. The
bootstrap 95\% CI for the five-year DFS risk difference of
levamisole plus 5-FU against observation is $-16.7$ percentage points
($-24.2$ to $-9.4$), corresponding to a number needed to treat of six
patients to prevent one additional DFS event at five years. Levamisole
alone shows no detectable absolute benefit ($-1.8$ points, $-9.3$ to
$+5.7$). These absolute magnitudes frame the trial's clinical
importance in a way the hazard ratio alone cannot.

\subsection{Proportional-hazards diagnostics}

The scaled Schoenfeld-residual analysis flags two covariates: tumour
differentiation ($p < 0.001$) and obstruction ($p = 0.008$), with a
significant global test ($\chi^2 = 42.2$, 14 df, $p < 0.001$).
Reassuringly, the treatment indicator itself shows no evidence of a
time-varying effect ($p = 0.51$), so the primary hazard-ratio summary
is defensible. Figure~\ref{fig:zph} shows the residual patterns; the
obstruction effect visibly weakens after year two, a common phenomenon
for early surgical complications whose prognostic impact fades. In
addition, replacing linear age and nodes with 4-degree-of-freedom
penalized splines significantly improves the log-likelihood
(likelihood-ratio $p < 0.001$), indicating threshold-like dose effects
in the nodal count that the linear term smooths over; the treatment
estimate is unaffected. Both diagnostics motivate the competing-risk
sensitivity analysis and justify keeping the primary estimand simple
while reporting these caveats.

\begin{figure}[htbp]
  \centering
  \includegraphics[width=0.92\textwidth, max height=0.36\textheight]{fig_zph.png}
  \caption{Proportional-hazards diagnostics: scaled Schoenfeld
  residuals with loess smooths for treatment, positive nodes, and
  obstruction (which exhibits a significant violation, $p = 0.008$).}
  \label{fig:zph}
\end{figure}

\subsection{Competing-risk sensitivity analysis}

Thirty-eight patients died without a recorded recurrence, and treating
these deaths as censoring---as the published coding does---formally
assumes they carry no recurrence-risk information. The competing-risk
view shows the practical consequence (Figure~\ref{fig:cif}): the
Aalen--Johansen cumulative incidence of recurrence sits below the naive
one-minus-Kaplan--Meier curve for every arm, with the gap reaching about
one percentage point at eight years in the observation arm.
The Fine--Gray subdistribution model, adjusted for the same covariates,
estimates the levamisole-plus-5-FU effect on the recurrence
subdistribution hazard at $0.596$ (95\% CI $0.469$--$0.756$)---
virtually identical to the cause-specific reading of $0.623$. The
conclusion of a substantial treatment benefit is therefore robust to
the endpoint-definition choice, and the analysis demonstrates the
magnitude of the naive estimator's upward bias in this dataset.

\begin{figure}[htbp]
  \centering
  \includegraphics[width=0.80\textwidth, max height=0.38\textheight]{fig_cif.png}
  \caption{Cumulative incidence of recurrence: Aalen--Johansen
  estimates treating death without recurrence as a competing event
  (solid) versus the naive one-minus-Kaplan--Meier estimator that
  censors such deaths (dashed).}
  \label{fig:cif}
\end{figure}
